% !TEX root = ../main.tex
\documentclass[../main.tex]{subfiles}

\begin{document}

\intermezzo{Triangulation and Cross-Modal Integration}

\begin{center}
\begin{tikzpicture}
    \fill[CoverGold] (-2,0) rectangle (-1.7,0.08);
    \fill[AccentPurple] (-1.5,0) rectangle (-1.2,0.08);
    \fill[AccentCyan] (-1,0) rectangle (-0.7,0.08);
    \fill[PandaBamboo] (-0.5,0) rectangle (-0.2,0.08);
    \fill[Azure] (0,0) rectangle (0.3,0.08);
\end{tikzpicture}
\end{center}

\vspace{0.5cm}

\begin{chapteroverview}
This chapter addresses the central methodological challenge of multimodal cartography: how do we integrate signals from five different data modalities into coherent role profiles? We develop a triangulation framework that treats convergence and divergence as equally informative, introduce quantitative integration methods, and demonstrate the approach through worked examples from the EIT ecosystem.
\end{chapteroverview}


%═══════════════════════════════════════════════════════════════════════════════
\section*{The Integration Challenge}
%═══════════════════════════════════════════════════════════════════════════════

Five modalities generate five ``views'' of each actor:

\sidenote{\textbf{Five Views:} \textcolor{Azure}{Websites} $\rightarrow$ Claimed, \textcolor{PandaBamboo}{News} $\rightarrow$ Attributed, \textcolor{Marigold}{Crunchbase} $\rightarrow$ Resourced, \textcolor{AccentCyan}{Social Media} $\rightarrow$ Relational, \textcolor{AccentPurple}{Visual} $\rightarrow$ Embodied.}

These views may \textbf{converge} (reinforcing confidence) or \textbf{diverge} (revealing tensions). Both patterns carry signal.


%═══════════════════════════════════════════════════════════════════════════════
\section*{Convergence as Validation}
%═══════════════════════════════════════════════════════════════════════════════

When multiple modalities agree, we gain confidence:

\begin{description}[leftmargin=0.5cm, itemsep=0.3em]
\item[Role Confirmation] An actor claiming ``connector'' on their website, positioned centrally in networks, and described as ``facilitator'' in news---this convergence validates the role.

\item[Cross-Modal Correlation] We compute correlation coefficients between role scores derived from different modalities. High correlation indicates robust measurement.

\item[Triangulation Strength] The more modalities that agree, the stronger the evidence for a role attribution.
\end{description}


%═══════════════════════════════════════════════════════════════════════════════
\section*{Divergence as Signal}
%═══════════════════════════════════════════════════════════════════════════════

When modalities disagree, we've found something interesting:

\begin{description}[leftmargin=0.5cm, itemsep=0.3em]
\item[Aspiration Gaps] Claimed roles (websites) exceed enacted roles (networks) = aspirational positioning

\item[Perception Gaps] Attributed roles (news) differ from claimed roles = legitimacy contests

\item[Resource Gaps] Claimed roles lack investment support = unfunded ambitions

\item[Strategic Gaps] Deliberate divergence for positioning purposes
\end{description}

\begin{insightbox}[The Connector Gap Revisited]
34\% of organizations claim connector roles on websites, but only 13\% occupy bridging network positions. This 21-point gap is not noise---it's data about intermediary aspiration exceeding intermediary reality.
\end{insightbox}


%═══════════════════════════════════════════════════════════════════════════════
\section*{Integration Methods}
%═══════════════════════════════════════════════════════════════════════════════

\subsection*{Weighted Aggregation}

Combine modality-specific role scores with weights reflecting reliability:

$$R_{\text{integrated}} = \sum_{m} w_m \cdot R_m$$

where $w_m$ weights each modality by data quality and theoretical relevance.

\subsection*{Profile Vectors}

Represent each actor as a vector across modalities and role dimensions:

$$\mathbf{v}_{\text{actor}} = [R_1^{\text{web}}, R_1^{\text{news}}, ..., R_k^{\text{visual}}]$$

Similarity between actors computed via cosine distance.

\subsection*{Gap Metrics}

Quantify divergence between specific modality pairs:

$$\text{Gap}_{ij} = |R_i - R_j|$$

High gaps flag actors for qualitative investigation.


%═══════════════════════════════════════════════════════════════════════════════
\section*{Chapter Summary}
%═══════════════════════════════════════════════════════════════════════════════

\begin{summarybox}
\begin{itemize}[leftmargin=*, itemsep=0.3em]
\item \textbf{Integration} is the central challenge of multimodal cartography.

\item \textbf{Convergence} across modalities validates role attributions and strengthens confidence.

\item \textbf{Divergence} reveals aspiration gaps, perception gaps, and strategic positioning.

\item \textbf{Methods} include weighted aggregation, profile vectors, and gap metrics.

\item \textbf{Both convergence and divergence} are informative---neither should be dismissed as noise.
\end{itemize}
\end{summarybox}

\vspace{0.5cm}

\begin{center}
\begin{tikzpicture}
    \fill[CoverGold] (-2,0) rectangle (-1.7,0.08);
    \fill[AccentPurple] (-1.5,0) rectangle (-1.2,0.08);
    \fill[AccentCyan] (-1,0) rectangle (-0.7,0.08);
    \fill[PandaBamboo] (-0.5,0) rectangle (-0.2,0.08);
    \fill[Azure] (0,0) rectangle (0.3,0.08);
\end{tikzpicture}
\end{center}

\closeintermezzo

\end{document}
